\documentclass[11pt]{article}
\usepackage[a4paper,margin=1.9cm,top=2.4cm,bottom=2.2cm]{geometry}
\usepackage{amsmath,amssymb,amsfonts,fontspec,microtype,fancyhdr,lastpage,enumitem,needspace}
\usepackage[most]{tcolorbox}
\usepackage{xcolor}
\definecolor{ink}{HTML}{1F3A5F}\definecolor{soft}{HTML}{EEF3F9}\definecolor{ans}{HTML}{F3F8F1}\definecolor{ansb}{HTML}{3C7A3C}
\pagestyle{fancy}\fancyhf{}
\lhead{\small\color{ink}\textbf{Linear Algebra I} \textperiodcentered\ Week 2 --- Matrices over $\mathbb{C}$, Conjugate Transpose, Abstract Spaces \& Subspace Tests}
\rhead{}\cfoot{\small Mr Malek Thiab \textperiodcentered\ Worksheet \textperiodcentered\ Page \thepage\ of \pageref{LastPage}}
\renewcommand{\headrulewidth}{0.4pt}
\setlength{\parindent}{0pt}\setlength{\parskip}{4pt}
\newcommand{\lv}[1]{\textsc{\small #1}}
\begin{document}
{\Large\bfseries\color{ink} Linear Algebra I --- Week 2}\\[2pt]
{\large\bfseries\color{ink} Matrices over $\mathbb{C}$, Conjugate Transpose, Abstract Spaces \& Subspace Tests}\\[3pt]
{\small Name: \underline{\hspace{5cm}}\hfill Date: \underline{\hspace{3cm}}\qquad All work \textbf{by hand} --- no calculator, no software.}
\vspace{4pt}\hrule\vspace{8pt}

\begin{tcolorbox}[colback=soft,colframe=ink,title={\bfseries Key Facts}]
\textbf{Complex numbers.} $i^{2}=-1$; $\overline{a+bi}=a-bi$; $z\overline z=|z|^{2}=a^{2}+b^{2}$; $\dfrac{z}{w}=\dfrac{z\,\overline w}{|w|^{2}}$. Also $\overline{z+w}=\overline z+\overline w$, $\overline{zw}=\overline z\,\overline w$, and $z=\overline z\iff z\in\mathbb{R}$.

\textbf{Matrices over $\mathbb{C}$.} $M_{m\times n}(\mathbb{C})$ is the set of $m\times n$ matrices with complex entries. Addition and scalar multiplication are entrywise; $(AB)_{jk}=\sum_{l}A_{jl}B_{lk}$. Multiplication is associative but \emph{not} commutative.

\textbf{Conjugate transpose.} $(A^{*})_{jk}=\overline{A_{kj}}$. Rules: $(A+B)^{*}=A^{*}+B^{*}$, $(\lambda A)^{*}=\overline{\lambda}A^{*}$, $(AB)^{*}=B^{*}A^{*}$, $(A^{*})^{*}=A$.
$A$ is \textbf{Hermitian} if $A^{*}=A$; \textbf{skew-Hermitian} if $A^{*}=-A$.

\textbf{Vector space over $\mathbb{F}$.} A set $V$ with $+$ and scalar multiplication satisfying: commutativity and associativity of $+$; a zero vector $\mathbf 0$; additive inverses; $1v=v$; $(\lambda\mu)v=\lambda(\mu v)$; $\lambda(u+v)=\lambda u+\lambda v$; $(\lambda+\mu)v=\lambda v+\mu v$.
Standard examples: $\mathbb{F}^{n}$, $M_{m\times n}(\mathbb{F})$, $P_{n}(\mathbb{F})$ (polynomials of degree $\le n$), $\mathcal F(S,\mathbb{F})$ (functions $S\to\mathbb{F}$, pointwise operations).

\textbf{Subspace test.} $W\subseteq V$ is a subspace iff (i) $\mathbf 0\in W$, (ii) $u,v\in W\Rightarrow u+v\in W$, (iii) $\lambda\in\mathbb{F},\,v\in W\Rightarrow\lambda v\in W$. To disprove: exhibit one failure. The field matters: a set can be a subspace over $\mathbb{R}$ but not over $\mathbb{C}$.
\end{tcolorbox}

\newpage
\section*{\large\color{ink} Worked Examples}
\begin{tcolorbox}[colback=white,colframe=ink!60,title={\bfseries Example 1 \textperiodcentered\ Complex arithmetic}]Compute $(2-i)(1+3i)$ and $\dfrac{5+i}{2-i}$.
\par\smallskip
\textbf{Solution.} \begin{enumerate}[leftmargin=1.6em,itemsep=1pt,topsep=2pt,label=\arabic*.]\item $(2-i)(1+3i)=2+6i-i-3i^{2}=5+5i$.
\item $\dfrac{5+i}{2-i}=\dfrac{(5+i)(2+i)}{(2-i)(2+i)}=\dfrac{10+5i+2i+i^{2}}{5}=\dfrac{9+7i}{5}$.
\end{enumerate}\textbf{Answer:} $5+5i$ and $\tfrac95+\tfrac75i$
\end{tcolorbox}
\begin{tcolorbox}[colback=white,colframe=ink!60,title={\bfseries Example 2 \textperiodcentered\ Product over $\mathbb{C}$}]Compute $\begin{pmatrix}i&1\\0&2\end{pmatrix}\begin{pmatrix}1&i\\i&0\end{pmatrix}$.
\par\smallskip
\textbf{Solution.} \begin{enumerate}[leftmargin=1.6em,itemsep=1pt,topsep=2pt,label=\arabic*.]\item Row 1: $(i\cdot1+1\cdot i,\ i\cdot i+1\cdot0)=(2i,\,-1)$.
\item Row 2: $(0\cdot1+2\cdot i,\ 0\cdot i+2\cdot0)=(2i,\,0)$.
\end{enumerate}\textbf{Answer:} $\begin{pmatrix}2i&-1\\2i&0\end{pmatrix}$
\end{tcolorbox}
\begin{tcolorbox}[colback=white,colframe=ink!60,title={\bfseries Example 3 \textperiodcentered\ Conjugate transpose}]Find $M^{*}$ for $M=\begin{pmatrix}2-i&3\\i&1+2i\end{pmatrix}$.
\par\smallskip
\textbf{Solution.} \begin{enumerate}[leftmargin=1.6em,itemsep=1pt,topsep=2pt,label=\arabic*.]\item Transpose: $\begin{pmatrix}2-i&i\\3&1+2i\end{pmatrix}$.
\item Conjugate every entry.
\end{enumerate}\textbf{Answer:} $M^{*}=\begin{pmatrix}2+i&-i\\3&1-2i\end{pmatrix}$
\end{tcolorbox}
\begin{tcolorbox}[colback=white,colframe=ink!60,title={\bfseries Example 4 \textperiodcentered\ Hermitian test}]Is $N=\begin{pmatrix}1&2+i\\2-i&3\end{pmatrix}$ Hermitian?
\par\smallskip
\textbf{Solution.} \begin{enumerate}[leftmargin=1.6em,itemsep=1pt,topsep=2pt,label=\arabic*.]\item $N^{*}$: transpose swaps $2+i$ and $2-i$; conjugating gives $\overline{2-i}=2+i$ in position $(1,2)$ and $\overline{2+i}=2-i$ in position $(2,1)$.
\item Diagonal entries $1,3$ are real and unchanged.
\end{enumerate}\textbf{Answer:} Yes, $N^{*}=N$.
\end{tcolorbox}
\begin{tcolorbox}[colback=white,colframe=ink!60,title={\bfseries Example 5 \textperiodcentered\ Proof: $(AB)^{*}=B^{*}A^{*}$}]Prove $(AB)^{*}=B^{*}A^{*}$ for $A\in M_{m\times n}(\mathbb{C})$, $B\in M_{n\times p}(\mathbb{C})$.
\par\smallskip
\textbf{Solution.} \begin{enumerate}[leftmargin=1.6em,itemsep=1pt,topsep=2pt,label=\arabic*.]\item Entry $(j,k)$ of the left side: $\overline{(AB)_{kj}}=\overline{\sum_{l}A_{kl}B_{lj}}=\sum_{l}\overline{A_{kl}}\;\overline{B_{lj}}$ (conjugation respects sums and products).
\item Entry $(j,k)$ of the right side: $\sum_{l}(B^{*})_{jl}(A^{*})_{lk}=\sum_{l}\overline{B_{lj}}\;\overline{A_{kl}}$.
\item The two sums agree term by term (multiplication in $\mathbb{C}$ commutes).
\end{enumerate}\textbf{Answer:} $(AB)^{*}=B^{*}A^{*}$.
\end{tcolorbox}
\begin{tcolorbox}[colback=white,colframe=ink!60,title={\bfseries Example 6 \textperiodcentered\ Axiom failure}]On $\mathbb{R}^{2}$ with usual $\oplus$ and $\lambda\odot(a,b)=(\lambda a,\,b)$: is it a vector space?
\par\smallskip
\textbf{Solution.} \begin{enumerate}[leftmargin=1.6em,itemsep=1pt,topsep=2pt,label=\arabic*.]\item Test $(\lambda+\mu)\odot v=\lambda\odot v\oplus\mu\odot v$ with $v=(a,b)$.
\item Left: $((\lambda+\mu)a,\ b)$. Right: $(\lambda a+\mu a,\ b+b)=((\lambda+\mu)a,\ 2b)$.
\item With $b\neq0$ these differ.
\end{enumerate}\textbf{Answer:} No: distributivity over scalar addition fails.
\end{tcolorbox}
\begin{tcolorbox}[colback=white,colframe=ink!60,title={\bfseries Example 7 \textperiodcentered\ Subspace test}]Show $W=\{p\in P_{2}(\mathbb{R}):p(0)=0\}$ is a subspace.
\par\smallskip
\textbf{Solution.} \begin{enumerate}[leftmargin=1.6em,itemsep=1pt,topsep=2pt,label=\arabic*.]\item $\mathbf 0(0)=0$, so $\mathbf 0\in W$.
\item $p(0)=q(0)=0\Rightarrow(p+q)(0)=0$.
\item $p(0)=0\Rightarrow(\lambda p)(0)=\lambda\cdot0=0$.
\end{enumerate}\textbf{Answer:} $W$ is a subspace.
\end{tcolorbox}
\begin{tcolorbox}[colback=white,colframe=ink!60,title={\bfseries Example 8 \textperiodcentered\ Non-subspace}]Is $D=\{A\in M_{2\times2}(\mathbb{R}):\det A=0\}$ a subspace?
\par\smallskip
\textbf{Solution.} \begin{enumerate}[leftmargin=1.6em,itemsep=1pt,topsep=2pt,label=\arabic*.]\item $A=\begin{pmatrix}1&0\\0&0\end{pmatrix}$ and $B=\begin{pmatrix}0&0\\0&1\end{pmatrix}$ both have $\det=0$, so lie in $D$.
\item $A+B=I$, and $\det I=1\neq0$.
\end{enumerate}\textbf{Answer:} No: not closed under addition.
\end{tcolorbox}
\newpage
\section*{\large\color{ink} Practice Problems}
\Needspace{10\baselineskip}\subsection*{\normalsize\color{ink} Part A --- Complex Arithmetic and Matrices over $\mathbb{C}$}
\Needspace{6\baselineskip}\begin{minipage}{\linewidth}\textbf{1.} \hfill{\scriptsize\color{ink}[Foundational]}\par\nopagebreak Compute $(3+2i)(1-4i)$ and write it in the form $a+bi$.
\vspace{2.2cm}\end{minipage}\par\vspace{4pt}
\Needspace{6\baselineskip}\begin{minipage}{\linewidth}\textbf{2.} \hfill{\scriptsize\color{ink}[Foundational]}\par\nopagebreak Compute $\dfrac{2+i}{1-3i}$ and write it in the form $a+bi$.
\vspace{2.6cm}\end{minipage}\par\vspace{4pt}
\Needspace{6\baselineskip}\begin{minipage}{\linewidth}\textbf{3.} \hfill{\scriptsize\color{ink}[Foundational]}\par\nopagebreak Let $A=\begin{pmatrix}1&i\\2-i&3\end{pmatrix}$ and $B=\begin{pmatrix}i&0\\1&1+i\end{pmatrix}$. Compute $A+B$ and $2i\,A$.
\vspace{3.4cm}\end{minipage}\par\vspace{4pt}
\Needspace{6\baselineskip}\begin{minipage}{\linewidth}\textbf{4.} \hfill{\scriptsize\color{ink}[Intermediate]}\par\nopagebreak With $A$ and $B$ as in Problem 3, compute $AB$.
\vspace{3.4cm}\end{minipage}\par\vspace{4pt}
\Needspace{6\baselineskip}\begin{minipage}{\linewidth}\textbf{5.} \hfill{\scriptsize\color{ink}[Foundational]}\par\nopagebreak Find the conjugate transpose $C^{*}$ of $C=\begin{pmatrix}1+i&2\\3i&4-i\\0&-2i\end{pmatrix}$, and state its size.
\vspace{3.4cm}\end{minipage}\par\vspace{4pt}
\Needspace{6\baselineskip}\begin{minipage}{\linewidth}\textbf{6.} \hfill{\scriptsize\color{ink}[Intermediate]}\par\nopagebreak Decide whether each matrix is Hermitian ($M^{*}=M$), skew-Hermitian ($M^{*}=-M$), or neither: \ $H=\begin{pmatrix}2&1-i\\1+i&5\end{pmatrix}$, \ $K=\begin{pmatrix}i&1\\-1&2i\end{pmatrix}$.
\vspace{4.0cm}\end{minipage}\par\vspace{4pt}
\Needspace{6\baselineskip}\begin{minipage}{\linewidth}\textbf{7.} \hfill{\scriptsize\color{ink}[Intermediate]}\par\nopagebreak Find real numbers $x,y$ such that $M=\begin{pmatrix}2&x+iy\\3-2i&5\end{pmatrix}$ is Hermitian.
\vspace{3.4cm}\end{minipage}\par\vspace{4pt}
\Needspace{6\baselineskip}\begin{minipage}{\linewidth}\textbf{8.} \hfill{\scriptsize\color{ink}[Challenge]}\par\nopagebreak Let $A=\begin{pmatrix}1&i\\2&1-i\end{pmatrix}$. Compute $\operatorname{tr}(A^{*}A)$ without computing all of $A^{*}A$, and explain why the result is always a non-negative real number.
\vspace{4.0cm}\end{minipage}\par\vspace{4pt}
\Needspace{10\baselineskip}\subsection*{\normalsize\color{ink} Part B --- Matrix Algebra over $\mathbb{C}$}
\Needspace{6\baselineskip}\begin{minipage}{\linewidth}\textbf{9.} \hfill{\scriptsize\color{ink}[Foundational]}\par\nopagebreak Let $A=\begin{pmatrix}1&i\\0&1\end{pmatrix}$. Compute $A^{2}$ and $A^{3}$, then state $A^{n}$.
\vspace{3.4cm}\end{minipage}\par\vspace{4pt}
\Needspace{6\baselineskip}\begin{minipage}{\linewidth}\textbf{10.} \hfill{\scriptsize\color{ink}[Intermediate]}\par\nopagebreak Let $A=\begin{pmatrix}0&i\\1&0\end{pmatrix}$ and $B=\begin{pmatrix}1&0\\0&-1\end{pmatrix}$. Compute $AB$, $BA$ and $AB-BA$.
\vspace{4.0cm}\end{minipage}\par\vspace{4pt}
\Needspace{6\baselineskip}\begin{minipage}{\linewidth}\textbf{11.} \hfill{\scriptsize\color{ink}[Intermediate]}\par\nopagebreak Let $A=\begin{pmatrix}1&i\\0&1\end{pmatrix}$ and $B=\begin{pmatrix}1&0\\i&1\end{pmatrix}$. Verify $(AB)^{*}=B^{*}A^{*}$ by computing both sides.
\vspace{4.4cm}\end{minipage}\par\vspace{4pt}
\Needspace{6\baselineskip}\begin{minipage}{\linewidth}\textbf{12.} \hfill{\scriptsize\color{ink}[Challenge]}\par\nopagebreak Let $A$ be an $n\times n$ Hermitian matrix. Prove that (a) every diagonal entry of $A$ is real, and (b) $iA$ is skew-Hermitian.
\vspace{4.6cm}\end{minipage}\par\vspace{4pt}
\Needspace{10\baselineskip}\subsection*{\normalsize\color{ink} Part C --- Abstract Vector Spaces}
\Needspace{6\baselineskip}\begin{minipage}{\linewidth}\textbf{13.} \hfill{\scriptsize\color{ink}[Intermediate]}\par\nopagebreak On $\mathbb{R}^{2}$ define $(a,b)\oplus(c,d)=(a+c,\,b+d)$ and $\lambda\odot(a,b)=(\lambda a,\,0)$. Is $(\mathbb{R}^{2},\oplus,\odot)$ a vector space over $\mathbb{R}$? Justify.
\vspace{3.4cm}\end{minipage}\par\vspace{4pt}
\Needspace{6\baselineskip}\begin{minipage}{\linewidth}\textbf{14.} \hfill{\scriptsize\color{ink}[Challenge]}\par\nopagebreak Let $V=\mathbb{R}_{>0}$ with $x\oplus y=xy$ and $\lambda\odot x=x^{\lambda}$. (Assume it is a real vector space.) Find the zero vector, the additive inverse of $5$, and compute $(3\odot2)\oplus\big((-1)\odot5\big)$.
\vspace{5.0cm}\end{minipage}\par\vspace{4pt}
\Needspace{6\baselineskip}\begin{minipage}{\linewidth}\textbf{15.} \hfill{\scriptsize\color{ink}[Challenge]}\par\nopagebreak Using only the vector space axioms, prove that $0\cdot v=\mathbf 0$ for every vector $v$ (here $0$ is the scalar zero and $\mathbf 0$ the zero vector).
\vspace{5.0cm}\end{minipage}\par\vspace{4pt}
\Needspace{6\baselineskip}\begin{minipage}{\linewidth}\textbf{16.} \hfill{\scriptsize\color{ink}[Foundational]}\par\nopagebreak In $P_{2}(\mathbb{R})$ let $p(x)=1+x-x^{2}$ and $q(x)=2-x+x^{2}$. Compute $3p-2q$.
\vspace{2.6cm}\end{minipage}\par\vspace{4pt}
\Needspace{6\baselineskip}\begin{minipage}{\linewidth}\textbf{17.} \hfill{\scriptsize\color{ink}[Foundational]}\par\nopagebreak In $\mathcal F(\mathbb{R},\mathbb{R})$ let $f(x)=x^{2}$ and $g(x)=3x+1$. Write the function $2f-g$ and evaluate it at $x=2$.
\vspace{2.6cm}\end{minipage}\par\vspace{4pt}
\Needspace{6\baselineskip}\begin{minipage}{\linewidth}\textbf{18.} \hfill{\scriptsize\color{ink}[Intermediate]}\par\nopagebreak In $\mathbb{C}^{2}$ over $\mathbb{C}$, compute $(1+i)\,(2-i,\;3i)+(-i)\,(1,\;1+i)$.
\vspace{3.4cm}\end{minipage}\par\vspace{4pt}
\Needspace{10\baselineskip}\subsection*{\normalsize\color{ink} Part D --- Subspace Tests}
\Needspace{6\baselineskip}\begin{minipage}{\linewidth}\textbf{19.} \hfill{\scriptsize\color{ink}[Foundational]}\par\nopagebreak Show that $W=\{(x,y,z)\in\mathbb{R}^{3}: x+2y-z=0\}$ is a subspace of $\mathbb{R}^{3}$.
\vspace{4.4cm}\end{minipage}\par\vspace{4pt}
\Needspace{6\baselineskip}\begin{minipage}{\linewidth}\textbf{20.} \hfill{\scriptsize\color{ink}[Foundational]}\par\nopagebreak Show that $W=\{(x,y)\in\mathbb{R}^{2}: xy=0\}$ is \emph{not} a subspace of $\mathbb{R}^{2}$.
\vspace{3.4cm}\end{minipage}\par\vspace{4pt}
\Needspace{6\baselineskip}\begin{minipage}{\linewidth}\textbf{21.} \hfill{\scriptsize\color{ink}[Challenge]}\par\nopagebreak Let $W=\{A\in M_{2\times2}(\mathbb{C}):A^{*}=A\}$ (Hermitian matrices). Is $W$ a subspace of $M_{2\times2}(\mathbb{C})$ over $\mathbb{C}$? Is it a subspace over $\mathbb{R}$? Justify both.
\vspace{5.0cm}\end{minipage}\par\vspace{4pt}
\Needspace{6\baselineskip}\begin{minipage}{\linewidth}\textbf{22.} \hfill{\scriptsize\color{ink}[Intermediate]}\par\nopagebreak In $P_{3}(\mathbb{R})$ let $W_1=\{p: p(2)=0\}$ and $W_2=\{p: p(2)=1\}$. Decide which is a subspace.
\vspace{4.4cm}\end{minipage}\par\vspace{4pt}
\Needspace{6\baselineskip}\begin{minipage}{\linewidth}\textbf{23.} \hfill{\scriptsize\color{ink}[Intermediate]}\par\nopagebreak Show that the set $E$ of even functions $f\in\mathcal F(\mathbb{R},\mathbb{R})$, i.e.\ $f(-x)=f(x)$ for all $x$, is a subspace.
\vspace{4.4cm}\end{minipage}\par\vspace{4pt}
\Needspace{6\baselineskip}\begin{minipage}{\linewidth}\textbf{24.} \hfill{\scriptsize\color{ink}[Challenge]}\par\nopagebreak In $\mathbb{R}^{2}$ let $U=\{(x,0)\}$ and $W=\{(0,y)\}$. Show that $U\cap W$ and $U+W$ are subspaces, but $U\cup W$ is not.
\vspace{5.0cm}\end{minipage}\par\vspace{4pt}
\end{document}
